\documentclass{../pset}

% User Information
\course{APS106H1: Fundamentals of Computer Programming}
\assignment{Problem Set 1}


\begin{document}

\begin{problem}
Write a program that simply prints "I am going to love this course" to the screen.
\end{problem}

\begin{problem}
Write a program that defines two integers, one being the largest 3-digit number and one being the smallest 5-digit number. Then, print these two values, and their sum to the
screen.
\end{problem}

\begin{problem}
Write a program that calculates the final cost of a set of items. For example, if each item costs 1.43 and we need 9 of them. Then, print the final cost.  

For example:
\begin{lstlisting}[language={}]
Cost? 1.43
How many? 9
Total cost is $12.87.
\end{lstlisting}


First do it by "hard-coding" the numbers. Once you have that working, add the ability to ask the user to input the numbers.
\end{problem}

\begin{problem}
Write a program that for a temperature in degrees Celsius, prints the equivalent temperature in degrees Fahrenheit (multiply C by 1.8, then add 32).  

For example:
\begin{lstlisting}[language={}]
Celsius temperature: 30
Fahrenheit equivalent is: 86
\end{lstlisting}

As with Problem 3, first put the Celsius temperature in your program directly and get it
working. Then add the functionality to prompt the user for the Celsius temperature.
\end{problem}

\begin{problem}
Write a program to evaluate $y = 4x + 3$ for any float number the user may pick. Ask a user to input any $x$, and then evaluate $y$, and present the result to two decimal places.

For example:
\begin{lstlisting}[language={}]
Input x: 3.69
The value of y = 4x + 3 is 17.76
\end{lstlisting}


\end{problem}


\end{document}